\phantomsection
\addcontentsline{toc}{chapter}{概论}
\markboth{概论}{概论}
\vspace*{1.2cm}
\begin{center}
{\LARGE\bfseries 概论\par}
\end{center}
\vspace{1.2em}

自由电子量子光学研究连续电子谱怎样同光场、材料激发和测量装置交换能量、动量与量子相干。要把这些过程放在同一套语言中，起点不能只是某一种边带模型或某一类器件：电子的运动受相对论时空对称与 Dirac 动力学约束，光和材料激发需要场量子化与 Green 张量，实验读出又离不开散射、开放系统和量子测量。讲义沿着这条物理链展开，从共同的数学工具和相对论场论进入 QED、QCD 与标准模型，再回到自由电子与经典场、量子光场、周期结构和介质的相互作用，最后讨论传播、辐射、测量与反演。随着具体尺度、几何和受控近似逐步加入，同一母理论会自然化成不同模型；这些模型最终都要落到实验真正能够读出的电子能谱、光子计数以及时空结构上。

QCD 和标准模型在这里承担方法训练的作用。QCD 的颜色平均、因子化与尺度演化说明，微观振幅必须经过初态统计和长距离矩阵元才能成为实验截面；弱衰变与中微子散射则展示手征投影、质量阈值和有限探测窗口怎样改变观测结果。\secref{sec:sm-neutrino-matter-applications} 进一步用中微子物质传播连接前向散射、矩阵传播子与绝热近似。回到电子光学时，同样需要把相互作用、传播和探测依次计算，只是强子 PDF 被电子入射态与材料响应所替代。

初次阅读可沿各节正文掌握对象定义、主方程及适用范围，再用局部推导核对自旋迹、积分测度和近似余项。用于课堂练习时，\secref{sec:qcd-momentum-evolution-applications} 从分裂核算出夸克与胶子动量份额，\secref{sec:qcd-drell-yan-applications} 从部分子振幅计算轻子对事件数；\secref{sec:sm-pseudoscalar-decay-applications} 和\secref{sec:sm-neutrino-electron-applications} 分别计算介子衰变与电子反冲谱，\secref{sec:electron-colored-noise-applications} 则把同样的误差控制落实到边带相干。例题中的教学参数明确标注，涉及精密实验比较时还须加入相应的辐射修正和探测响应。

自由电子部分可沿一条可复算的实验链阅读。\secref{sec:electron-channel-coherence-worked} 先用波包重叠和偏迹区分能谱概率与量子相干，\secref{sec:electron-finite-window-worked} 再计算有限作用时间造成的谱宽及黄金律误差。\secref{sec:electron-pinem-field-calibration} 把实场峰值换成耦合与边带概率，\secref{sec:electron-quantum-exchange-coherence} 计算量子光场使电子纯度下降的幅度。随后，\secref{sec:electron-recoil-bragg-calculation} 对照完整反冲格与两态近似，\secref{sec:electron-temporal-bunching-worked} 从能谱相干算出脉冲列；\secref{sec:electron-reference-tomography-worked} 和\secref{sec:electron-detector-inversion-worked} 最后把四相位测量、有限分辨率和有限收集窗口接到电子相干与近场复振幅的反演上。


本书例题以保留符号参数的解析计算为主：先确定待求对象，再展示配方积分、矩阵指数、生成函数或误差不等式的关键步骤，最后检查特殊极限。\secref{sec:electron-thermal-exchange-analytic} 用热态偏迹与生成函数求完整电子交换分布；\secref{sec:electron-gaussian-lens-analytic} 从连续 TDSE 解出有限波包的聚焦与不确定关系。数值传播图用于检验解析近似何时成立，不要求读者以数值运算代替例题推导。
